% !TEX root = ../../../main.tex
% 来源：中学数学实验教材·第二册（几何）
% 章节：实验几何 / 方向、角度与平行
% 原始文件：1.tex，行 999-1024
% 类型：tikz
% ---
\begin{tikzpicture}[scale=2]
\begin{scope}
\draw(-1,.5)node[left]{$A$}--(0,0)node[right]{$O$}--(-1,-.5)node[left]{$B$};
\node at (-.5,-1){(1)};
\end{scope}
\begin{scope}[xshift=1.3cm]
	\tkzDefPoints{0/.5/O, -1/-.5/A, -.5/-.7/B, .5/-.7/C, 1/-.5/D}
	\tkzDrawSegments(O,A O,B O,C O,D)
	\tkzMarkAngles[mark=none, size=.4](A,O,B C,O,D)
	\tkzLabelAngle[pos=.5](A,O,B){1}
	\tkzLabelAngle[pos=.5](B,O,C){2}
	\tkzLabelAngle[pos=.5](C,O,D){3}
	\tkzMarkAngles[mark=none, size=.3](B,O,C)
	\node at (0,-1){(2)};
\end{scope}
\begin{scope}[xshift=3.5cm]
	\tkzDefPoints{0/0/O, -1/.5/A, 1/-.5/B, 1/.5/C, -1/-.5/D}
	\tkzDrawSegments(A,B C,D)
	\tkzMarkAngles[mark=none, size=.2](C,O,A D,O,B)
	\tkzMarkAngles[mark=none, size=.3](A,O,D)
	\tkzLabelAngle[pos=.3](D,O,B){$\gamma$}
	\tkzLabelAngle[pos=.4](A,O,D){$\beta$}
	\tkzLabelAngle[pos=.3](C,O,A){$\alpha$}
	\node at (0,-1){(3)};
\end{scope}
\end{tikzpicture}
